\documentclass[11pt]{article}
\usepackage[letterpaper,top=0.68in,bottom=0.68in,left=0.78in,right=0.78in]{geometry}
% Shared style for Advisor_Packet documents.
% Input after \documentclass and geometry (or let documents set geometry first).
% Documents should: \usepackage{...} then % Shared style for Advisor_Packet documents.
% Input after \documentclass and geometry (or let documents set geometry first).
% Documents should: \usepackage{...} then % Shared style for Advisor_Packet documents.
% Input after \documentclass and geometry (or let documents set geometry first).
% Documents should: \usepackage{...} then \input{../Assets/packet_style.tex}

\usepackage[T1]{fontenc}
\usepackage[utf8]{inputenc}
\usepackage{lmodern}
\usepackage{microtype}
\usepackage{amsmath,amssymb}
\usepackage{booktabs,tabularx,array,colortbl}
\usepackage{xcolor}
\usepackage{enumitem}
\usepackage{hyperref}
\usepackage{titlesec}
\usepackage{fancyhdr}
\usepackage{ragged2e}
\usepackage{needspace}
\usepackage[most]{tcolorbox}

\definecolor{navy}{HTML}{1B3A5C}
\definecolor{slate}{HTML}{5B6B7C}
\definecolor{bluegray}{HTML}{E4EBF3}
\definecolor{lightgray}{HTML}{F4F5F7}
\definecolor{accent}{HTML}{6B7F93}
\definecolor{panelgreen}{HTML}{E8F0EA}
\definecolor{panelborder}{HTML}{C5D0DB}

\hypersetup{colorlinks=true,linkcolor=navy,urlcolor=navy}
\setlist[itemize]{leftmargin=1.35em,itemsep=2pt,topsep=2pt,parsep=0pt}
\setlist[enumerate]{leftmargin=1.55em,itemsep=2pt,topsep=2pt,parsep=0pt}
\setlength{\parindent}{0pt}
\setlength{\parskip}{3.5pt}
\setlength{\headheight}{22pt}
\raggedbottom

\titleformat{\section}{\large\bfseries\color{navy}}{}{0pt}{}[\vspace{1pt}{\color{panelborder}\titlerule[0.6pt]}]
\titleformat{\subsection}{\normalsize\bfseries\color{navy}}{}{0pt}{}
\titlespacing*{\section}{0pt}{10pt}{5pt}
\titlespacing*{\subsection}{0pt}{7pt}{3pt}

\newcommand{\PacketHeader}[1]{%
  \pagestyle{fancy}%
  \fancyhf{}%
  \lhead{\small\color{slate} #1}%
  \rhead{\small\color{slate} Nicholas Mino}%
  \cfoot{\small\color{slate}\thepage}%
  \renewcommand{\headrulewidth}{0.4pt}%
  \renewcommand{\headrule}{\hbox to\headwidth{\color{panelborder}\leaders\hrule height \headrulewidth\hfill}}%
}

\tcbset{
  basebox/.style={
    enhanced, breakable,
    colback=lightgray, colframe=panelborder,
    boxrule=0.5pt, arc=1.5pt,
    left=8pt, right=8pt, top=5pt, bottom=5pt
  },
  notebox/.style={
    enhanced, breakable=false,
    colback=bluegray, colframe=navy,
    boxrule=0.55pt, arc=1.5pt,
    left=8pt, right=8pt, top=5pt, bottom=5pt
  },
  prioritybox/.style={
    enhanced, breakable=false,
    colback=bluegray, colframe=navy,
    boxrule=0.9pt, arc=1.5pt,
    left=8pt, right=8pt, top=5pt, bottom=5pt,
    fonttitle=\bfseries\color{navy},
    colbacktitle=bluegray, coltitle=navy,
    attach boxed title to top left={yshift=-1.5pt, xshift=6pt},
    boxed title style={boxrule=0pt, arc=1pt, left=5pt, right=5pt, top=1.5pt, bottom=1.5pt}
  },
  entrybox/.style={
    enhanced, breakable,
    colback=white, colframe=panelborder,
    boxrule=0.55pt, arc=1.5pt,
    left=7pt, right=7pt, top=4pt, bottom=4pt,
    before skip=5pt, after skip=3pt,
    fonttitle=\bfseries\color{navy}\normalsize,
    colbacktitle=bluegray, coltitle=navy,
    attach boxed title to top left={yshift=-2pt, xshift=7pt},
    boxed title style={boxrule=0pt, arc=1pt, left=5pt, right=5pt, top=1.5pt, bottom=1.5pt}
  }
}


\usepackage[T1]{fontenc}
\usepackage[utf8]{inputenc}
\usepackage{lmodern}
\usepackage{microtype}
\usepackage{amsmath,amssymb}
\usepackage{booktabs,tabularx,array,colortbl}
\usepackage{xcolor}
\usepackage{enumitem}
\usepackage{hyperref}
\usepackage{titlesec}
\usepackage{fancyhdr}
\usepackage{ragged2e}
\usepackage{needspace}
\usepackage[most]{tcolorbox}

\definecolor{navy}{HTML}{1B3A5C}
\definecolor{slate}{HTML}{5B6B7C}
\definecolor{bluegray}{HTML}{E4EBF3}
\definecolor{lightgray}{HTML}{F4F5F7}
\definecolor{accent}{HTML}{6B7F93}
\definecolor{panelgreen}{HTML}{E8F0EA}
\definecolor{panelborder}{HTML}{C5D0DB}

\hypersetup{colorlinks=true,linkcolor=navy,urlcolor=navy}
\setlist[itemize]{leftmargin=1.35em,itemsep=2pt,topsep=2pt,parsep=0pt}
\setlist[enumerate]{leftmargin=1.55em,itemsep=2pt,topsep=2pt,parsep=0pt}
\setlength{\parindent}{0pt}
\setlength{\parskip}{3.5pt}
\setlength{\headheight}{22pt}
\raggedbottom

\titleformat{\section}{\large\bfseries\color{navy}}{}{0pt}{}[\vspace{1pt}{\color{panelborder}\titlerule[0.6pt]}]
\titleformat{\subsection}{\normalsize\bfseries\color{navy}}{}{0pt}{}
\titlespacing*{\section}{0pt}{10pt}{5pt}
\titlespacing*{\subsection}{0pt}{7pt}{3pt}

\newcommand{\PacketHeader}[1]{%
  \pagestyle{fancy}%
  \fancyhf{}%
  \lhead{\small\color{slate} #1}%
  \rhead{\small\color{slate} Nicholas Mino}%
  \cfoot{\small\color{slate}\thepage}%
  \renewcommand{\headrulewidth}{0.4pt}%
  \renewcommand{\headrule}{\hbox to\headwidth{\color{panelborder}\leaders\hrule height \headrulewidth\hfill}}%
}

\tcbset{
  basebox/.style={
    enhanced, breakable,
    colback=lightgray, colframe=panelborder,
    boxrule=0.5pt, arc=1.5pt,
    left=8pt, right=8pt, top=5pt, bottom=5pt
  },
  notebox/.style={
    enhanced, breakable=false,
    colback=bluegray, colframe=navy,
    boxrule=0.55pt, arc=1.5pt,
    left=8pt, right=8pt, top=5pt, bottom=5pt
  },
  prioritybox/.style={
    enhanced, breakable=false,
    colback=bluegray, colframe=navy,
    boxrule=0.9pt, arc=1.5pt,
    left=8pt, right=8pt, top=5pt, bottom=5pt,
    fonttitle=\bfseries\color{navy},
    colbacktitle=bluegray, coltitle=navy,
    attach boxed title to top left={yshift=-1.5pt, xshift=6pt},
    boxed title style={boxrule=0pt, arc=1pt, left=5pt, right=5pt, top=1.5pt, bottom=1.5pt}
  },
  entrybox/.style={
    enhanced, breakable,
    colback=white, colframe=panelborder,
    boxrule=0.55pt, arc=1.5pt,
    left=7pt, right=7pt, top=4pt, bottom=4pt,
    before skip=5pt, after skip=3pt,
    fonttitle=\bfseries\color{navy}\normalsize,
    colbacktitle=bluegray, coltitle=navy,
    attach boxed title to top left={yshift=-2pt, xshift=7pt},
    boxed title style={boxrule=0pt, arc=1pt, left=5pt, right=5pt, top=1.5pt, bottom=1.5pt}
  }
}


\usepackage[T1]{fontenc}
\usepackage[utf8]{inputenc}
\usepackage{lmodern}
\usepackage{microtype}
\usepackage{amsmath,amssymb}
\usepackage{booktabs,tabularx,array,colortbl}
\usepackage{xcolor}
\usepackage{enumitem}
\usepackage{hyperref}
\usepackage{titlesec}
\usepackage{fancyhdr}
\usepackage{ragged2e}
\usepackage{needspace}
\usepackage[most]{tcolorbox}

\definecolor{navy}{HTML}{1B3A5C}
\definecolor{slate}{HTML}{5B6B7C}
\definecolor{bluegray}{HTML}{E4EBF3}
\definecolor{lightgray}{HTML}{F4F5F7}
\definecolor{accent}{HTML}{6B7F93}
\definecolor{panelgreen}{HTML}{E8F0EA}
\definecolor{panelborder}{HTML}{C5D0DB}

\hypersetup{colorlinks=true,linkcolor=navy,urlcolor=navy}
\setlist[itemize]{leftmargin=1.35em,itemsep=2pt,topsep=2pt,parsep=0pt}
\setlist[enumerate]{leftmargin=1.55em,itemsep=2pt,topsep=2pt,parsep=0pt}
\setlength{\parindent}{0pt}
\setlength{\parskip}{3.5pt}
\setlength{\headheight}{22pt}
\raggedbottom

\titleformat{\section}{\large\bfseries\color{navy}}{}{0pt}{}[\vspace{1pt}{\color{panelborder}\titlerule[0.6pt]}]
\titleformat{\subsection}{\normalsize\bfseries\color{navy}}{}{0pt}{}
\titlespacing*{\section}{0pt}{10pt}{5pt}
\titlespacing*{\subsection}{0pt}{7pt}{3pt}

\newcommand{\PacketHeader}[1]{%
  \pagestyle{fancy}%
  \fancyhf{}%
  \lhead{\small\color{slate} #1}%
  \rhead{\small\color{slate} Nicholas Mino}%
  \cfoot{\small\color{slate}\thepage}%
  \renewcommand{\headrulewidth}{0.4pt}%
  \renewcommand{\headrule}{\hbox to\headwidth{\color{panelborder}\leaders\hrule height \headrulewidth\hfill}}%
}

\tcbset{
  basebox/.style={
    enhanced, breakable,
    colback=lightgray, colframe=panelborder,
    boxrule=0.5pt, arc=1.5pt,
    left=8pt, right=8pt, top=5pt, bottom=5pt
  },
  notebox/.style={
    enhanced, breakable=false,
    colback=bluegray, colframe=navy,
    boxrule=0.55pt, arc=1.5pt,
    left=8pt, right=8pt, top=5pt, bottom=5pt
  },
  prioritybox/.style={
    enhanced, breakable=false,
    colback=bluegray, colframe=navy,
    boxrule=0.9pt, arc=1.5pt,
    left=8pt, right=8pt, top=5pt, bottom=5pt,
    fonttitle=\bfseries\color{navy},
    colbacktitle=bluegray, coltitle=navy,
    attach boxed title to top left={yshift=-1.5pt, xshift=6pt},
    boxed title style={boxrule=0pt, arc=1pt, left=5pt, right=5pt, top=1.5pt, bottom=1.5pt}
  },
  entrybox/.style={
    enhanced, breakable,
    colback=white, colframe=panelborder,
    boxrule=0.55pt, arc=1.5pt,
    left=7pt, right=7pt, top=4pt, bottom=4pt,
    before skip=5pt, after skip=3pt,
    fonttitle=\bfseries\color{navy}\normalsize,
    colbacktitle=bluegray, coltitle=navy,
    attach boxed title to top left={yshift=-2pt, xshift=7pt},
    boxed title style={boxrule=0pt, arc=1pt, left=5pt, right=5pt, top=1.5pt, bottom=1.5pt}
  }
}

\usepackage{xurl}
\PacketHeader{SURA Stabilized-Algorithm Catalog Overview}

\begin{document}

\begin{center}
{\LARGE\bfseries\color{navy} Stabilized-Algorithm Catalog Overview}\\[3pt]
{\normalsize\color{slate} Intellectual contribution of the literature extraction}\\[2pt]
{\small\textbf{Nick Mino}\quad\textcolor{accent}{|}\quad 29 July 2026}
\end{center}

\vspace{2pt}

\begin{tcolorbox}[notebox]
This overview summarizes the intellectual contribution of the catalog.
It does \emph{not} reproduce the catalog itself.
\end{tcolorbox}

\section{Purpose}

The catalog was built to identify recurring structures in how data-dependent
algorithms and selection procedures are stabilized---so that later theoretical work
could search for reusable mechanisms rather than inventing a one-off construction.
In the intellectual timeline it is recorded as a primary research product
(approximately 20 reconstructed hours within the June--July effort), not a minor
reading list.

\section{What was extracted}

Comparable fields used when processing sources included (where available):

\begin{itemize}
  \item source and citation;
  \item algorithm or selection rule;
  \item selected object;
  \item stability notion;
  \item assumptions;
  \item stabilization mechanism;
  \item noise, perturbation, randomization, or regularization;
  \item primitive decomposition;
  \item proof strategy;
  \item composition rule;
  \item utility cost;
  \item downstream inference or generalization guarantee;
  \item limitations;
  \item relevance to the current project.
\end{itemize}

\section{Representative entries (by mechanism)}

Entries below are chosen to illustrate \emph{distinct stabilization mechanisms}.
They are grounded in repository digests, the H1 catalog note, and the paper index.
They are not claimed to be exhaustive, and they are not themselves proofs of
Formulation~II.

\begin{tcolorbox}[entrybox, title=Randomized selection / noisy max]
\textbf{Paper / method.} Zrnic \& Jordan, \emph{Post-Selection Inference via Algorithmic Stability}
(digest: \texttt{Formulation/research/literature/digest\_zrnic\_jordan\_stability.md}).\\
\textbf{Selected object.} A (possibly randomized) selection map $\hat S$ (e.g.\ noisy /
stable screening or Stable LASSO support).\\
\textbf{Mechanism.} Explicit randomization (Laplace / report-noisy-max style; noisy
Frank--Wolfe) inducing $(\eta,\tau,\nu)$-stability of the selection law.\\
\textbf{Guarantee.} Stability transfers to post-selection CI validity with
$e^\eta$-adjusted classical level (Thm.~2 in the digest).\\
\textbf{Relevance.} Closest technical template for ``stabilize selection, then retain a
classical inferential object''---subject to advisor judgment on baselines.
\end{tcolorbox}

\begin{tcolorbox}[entrybox, title=Algorithmic stability (distributional)]
\textbf{Paper / method.} Same Zrnic--Jordan line; Bassily--Freund typical stability cited
as related (digest §1).\\
\textbf{Selected object.} Randomized algorithm $A$ outputting a selection.\\
\textbf{Mechanism.} Indistinguishability of the selection law from an oracle law on a
high-probability set (Def.~2).\\
\textbf{Guarantee.} Composition / post-processing closure; transfer lemmas to
post-selection coverage corrections.\\
\textbf{Relevance.} Clarifies that ``stability'' here is often about selection
\emph{laws}, which must be carefully related to Part~I replace-one movement of
$\hat\theta$.
\end{tcolorbox}

\begin{tcolorbox}[entrybox, title=Regularization / constrained randomized optimization]
\textbf{Paper / method.} Stable LASSO / stable marginal screening constructions in
Zrnic--Jordan (digest §2, design propositions).\\
\textbf{Selected object.} Model support or screened coordinates.\\
\textbf{Mechanism.} Optimization with injected noise and regularization-style control of
sensitivity; utility propositions quantify selection-quality cost.\\
\textbf{Guarantee.} Explicit $(\eta,\tau,\nu)$-stability rates with stated assumptions.\\
\textbf{Relevance.} Concrete precedent for randomizing a constrained program---a
structural cousin of randomizing problem~(2) in the write-up.
\end{tcolorbox}

\begin{tcolorbox}[entrybox, title=Conformal risk control (fixed hyperparameter)]
\textbf{Paper / method.} Angelopoulos et al., \emph{Conformal Risk Control}
(digest: \texttt{Formulation/research/literature/digest\_angelopoulos\_crc.md}).\\
\textbf{Selected object.} Threshold / risk-control parameter $\lambda$ when treated as
\emph{prespecified} (or chosen by the CRC rule under exchangeability).\\
\textbf{Mechanism.} Monotone risk + exchangeability; not a post-selection stabilizer
for a shared-data optimizer.\\
\textbf{Guarantee.} Fixed-parameter risk control / split-conformal coverage templates.\\
\textbf{Relevance.} Supplies the fixed-$\theta$ validity template behind write-up
assumption~(1); selecting $\hat\theta(\mathcal{D})$ on the same $\mathcal{D}$ breaks the
prespecified premise.
\end{tcolorbox}

\begin{tcolorbox}[entrybox, title=Differential privacy / max-information style control]
\textbf{Paper / method.} Max-divergence / $(\eta,\tau)$-indistinguishability in
Zrnic--Jordan Def.~1; H1 catalog note lists DP top-$k$ / Thresholdout-style examples among
MAIN catalog themes
(\texttt{Formulation/context/overview/swarm/04\_H1\_catalog.md};
overview literature integration).\\
\textbf{Selected object.} Private or information-bounded released selection.\\
\textbf{Mechanism.} Noise calibrated to sensitivity; privacy/max-info budgets as
stability parameters.\\
\textbf{Guarantee.} Distributional closeness of released outputs; often used as a bridge
to inference corrections.\\
\textbf{Relevance.} Suggests reusable accounting language ($\eta$, $\tau$, $\nu$) for
Part~II transfer statements---without identifying DP with the live UQ theorem.
\end{tcolorbox}

\begin{tcolorbox}[entrybox, title=Optimization robustness / structured selection noise]
\textbf{Paper / method.} Pattern notes synthesizing inverse-opt geometry with selection
randomization
(\texttt{Formulation/research/literature/patterns\_inverse\_opt\_structured\_noise.md});
lab CREDO/CREAM papers remain \emph{context}, not the live problem
(\texttt{Formulation/context/papers/INDEX.md}).\\
\textbf{Selected object.} Hyperparameter / decision variable from a constrained program.\\
\textbf{Mechanism.} Structured (not isotropic) noise on objectives/constraints, or soft
near-argmin sampling---aimed at selection stability.\\
\textbf{Guarantee.} Currently a \emph{design hypothesis} for randomizing write-up
problem~(2), not a completed theorem in this repository.\\
\textbf{Relevance.} Motivates Part~I assumptions and Part~II goal~(iii): randomize
selection; do not recalibrate $\mathcal{C}$.
\end{tcolorbox}

\section{Recurring patterns}

\paragraph{Supported by the extracted literature (observations).}
\begin{itemize}
  \item Margins / gaps often control selection instability (near-ties are hard).
  \item Strong convexity, regularization, or explicit noise can reduce optimizer /
        selector sensitivity.
  \item Randomized selection smooths discontinuous hard decisions (e.g.\ $\arg\max$).
  \item Stability parameters can transfer into inference corrections (coverage inflation,
        PoSI constants, etc.).
  \item Utility / selection quality is typically traded against inferential validity.
\end{itemize}

\paragraph{Hypotheses motivating the broader framework (not established theorems here).}
\begin{itemize}
  \item Primitive decompositions may support reusable proof modules across algorithms.
  \item Composition of primitives requires care; certificates need not be freely
        interchangeable.
  \item Structured noise tied to the geometry of constrained optimization may be more
        efficient than isotropic noise for stabilizing $\hat\theta$.
\end{itemize}

\section{Relation to the current formulation}

The catalog helped motivate Formulation~I and the later narrowing to Formulation~II.
It does \textbf{not} prove a universal stability framework, and operator-library results
do \textbf{not} automatically imply the Part~I--II UQ theorems.
The immediate research target is the narrower Part~I replace-one stability theorem under
explicit assumptions (see \texttt{Problem\_Writeup.pdf} and
\texttt{Research\_Update\_Summary.pdf}).

\section{Where to find the underlying material}

\begin{itemize}[itemsep=2pt]
  \item \path{Formulation/context/overview/swarm/04_H1_catalog.md} ---
        H1 working-survey summary ($\approx$25 MAIN examples claimed).
  \item \path{Formulation/context/overview/Research_Scope_Timeline_and_Authority.tex} ---
        Catalog status, primitives list, literature integration.
  \item \path{Formulation/research/literature/} ---
        Digests and pattern cards used for extraction.
  \item \path{Formulation/context/papers/INDEX.md} ---
        Prioritized paper corpus in the repository.
  \item \path{Advisor_Packet/Intellectual_Timeline/Intellectual_Timeline.pdf} ---
        Effort and role of the catalog in the research arc.
\end{itemize}

\vspace{4pt}
\noindent\textbf{Missing artifact.}
No full structured catalog spreadsheet/database dump was found in the repository on
29 July 2026. If a separate Overleaf/database export exists outside this tree, it is not
currently checked in; the packet therefore points to the surviving in-repo sources above.

\end{document}
